\documentclass{article}
\usepackage[utf8]{inputenc}
\usepackage[T1]{fontenc}
\usepackage{geometry}
\usepackage{graphicx}
\usepackage{float}
\usepackage{hyperref}
\usepackage[english,indonesian]{babel}

\title{Optimasi Pemrosesan Template LaTeX Menggunakan Pendekatan Berbasis AI}
\author{Academic Support Agent}
\date{\today}

\begin{document}
\maketitle

\begin{abstract}

\end{abstract}

\chapter{PENDAHULUAN}


\chapter{TINJAUAN PUSTAKA}


% Original content (commented out):
% \documentclass[a4paper, 12pt]{article}
\usepackage[T1]{fontenc} % T1 Font Encoding
\usepackage[utf8]{inputenc} % Input Encoding
\usepackage{tikz}
% \usepackage{todonotes}
% \setuptodonotes{inline}
\usetikzlibrary{shapes.geometric, arrows.meta, positioning}
\usepackage{amsmath}
\usepackage{listings}
\usepackage{adjustbox}

\usepackage{lineno}
% \linenumbers


\usepackage{lagadeild} % Include the custom style file
\addbibresource{refs.bib} % Add the bibliography resource
\addbibresource{domar.bib}

\def\thesisauthor{Eiginnafn Kenninafn} % Nafn
\def\thesistitle{Titill} % Titill
\def\thesissubtitle{\quad} % Undirtitill
\def\thesiskind{BA} % T.d. BA 
\def\thesissubject{lögfræði} % Fag
\def\department{Lagadeild} % Heiti deildar
\def\faculty{Félagsvísindasvið} % Heiti sviðs
\def\thesissupervisor{Nafn leiðbeinanda} % Nafn leiðbeinanda
\def\thesissupervisortitle{lögmaður} % titill leiðbeinenda
\def\thesisunits{6} % Fjöldi eininga
\def\month{Mánuður} % Mánuður
\def\year{Ár} % Ár


\begin{document}

\maketitlepage % Create the title page
\makesecondpage
\makethirdpage

\tableofcontents % Insert the table of contents
\thispagestyle{empty}
\newpage

\setcounter{page}{2}

\section{Inngangur}
Þetta sniðmát er fyrir BA ritgerðir við lagadeild Háskóla Íslands. Það er sniðið samkvæmt reglum lagadeildar eins og þær liggja fyrir í apríl 2024.

Gáttir allar, og það er ekki aðeins dyrnar sem hér eru til umræðu, því þær vísa í hugmyndir um það sem við göngum í gegnum þegar við leitum skilnings áður en fram er gengið.

\section{Meginmál}
Það er hægt að vitna í bækur, íslenskar og erlendar.\footcitebook[40-44]{sigurdhsson2015samningarettur}\textsuperscript{,}\footcitebook[12]{radin2013boilerplate} Það er líka hægt að vitna í útgefnar greinar.\footcitearticle[3]{bakos2009does} Einnig doktorsritgerðir.\footcitearticle[35-36]{kolbeinsson2021deep} Þessar heimildir eru geymdar í \texttt{refs.bib}.
Það er hægt að vitna í dóma: t.d. \textit{\cite{35-1981}} eða \textit{\cite{C-311/18}}. Dómar eru geymdir í \texttt{domar.bib}.

Elds er þörf, en ekki aðeins efnislegs elds, heldur líka þess andlega, þegar sá sem inn er kominn, og hér er ekki eingöngu átt við að komast inn í hús, heldur einnig að stíga inn í nýjan veruleika, á kné kalinn.

\subsection{Undirflokkur}
Undirflokkar bjóða upp á aukin sveigjanleika í uppsetningu.\footnotetext{Texti sem á heima í neðanmálsgrein.}
Vatns er þörf, en ekki aðeins til líkamlegrar endurnæringar, heldur líka til að hreinsa hugann þegar maður kemur til verðar, sem ekki aðeins er máltíð, heldur líka augnablik þar sem maður kemur saman til að ræða heimspekileg efni, eða til að freista þess að skilja meira.


\subsubsection{Undirundirflokkur}
Uppsetningin heldur áfram í undirundirflokkum. En hvað með þurru og þjóðlaðar, þar sem þessi orð geta haft dýpri merkingu, því að þau benda á þörfina fyrir það.

\section{Lokaorð}
Að lokum kemur eitthvað mikilvægt hér. Vits er þörf, því án þess væri það nær ómögulegt að rata, en víða ratar sá sem ekki þekkir sitt heimahérað eins og lófa sinn, því það er ekkert sem er eins auðvelt og það að sitja heima.

\newpage
\nolinenumbers
\printbibliography[heading=bibintoc, notkeyword=courtcase, notkeyword=eu, title={HEIMILDASKRÁ}]
\addcontentsline{toc}{section}{Heimildaskrá}  % Adjust the level if needed (e.g., chapter)

\newpage
% For the court cases
\printbibliography[keyword=courtcase, heading=bibcourt, env=italicbib]
\subsection*{\normalfont Dómar Evrópudómstólsins:}\vspace{-1em}
\printbibliography[keyword=eu, heading=empty, env=italicbib]
\addcontentsline{toc}{section}{Dómaskrá}

\end{document}




\end{document}